\documentclass[12pt]{article}
\usepackage[utf8]{inputenc}
\usepackage[T1]{fontenc}
\usepackage{hyperref}
\title{Tanjong Language Resources for K-12 Curriculum Development}
\author{Compiled by Claude Sonnet 5 (claude-sonnet-5) for Andr\'as Kornai}
\date{}
\begin{document}
\maketitle

\section{Introduction}
The code tnj (Glottocode tanj1248) is a Glottolog ``Bookkeeping'' entry, not a language in active use. Glottolog notes that tnj was retired in ISO 639-3 on 2008-01-14 and merged into Kanowit [kxn]. The merge was supported by 82\% lexical similarity between Kanowit and Tanjong. The surviving code is kxn, Kanowit-Tanjong Melanau. The prompt gave no family, region or speaker information. I did not search further for Kanowit or Tanjong Melanau materials, so no resources are listed and none are invented. Do not confuse tnj with tjn (Tonjon, an extinct Mande language of Cote d'Ivoire).

\section{Primary Educational Resources}
None searched or found.

\section{Linguistic and Reference Materials}
\begin{itemize}
\item Glottolog Bookkeeping entry for Tanjong: \url{https://glottolog.org/resource/languoid/id/tanj1248}
\end{itemize}

\section{Religious and Community Resources}
None found.

\section{Digital and Audio Resources}
None found.

\section{Curriculum Development Context}
Because tnj is a retired code, a separate curriculum for it is not meaningful. Work should be directed to Kanowit (kxn) instead, and I recommend that this entry be reviewed.

\end{document}